\input{preamble}

\begin{document}

\title{Probability}
\maketitle

\phantomsection
\label{section-phantom}

\tableofcontents

\section{Introduction}
\label{section-introduction}

\noindent
These notes develop probability theory.

\section{Robust and Efficient Mean
Estimation via Soft Shrinkage}
\label{section-robust-efficient-mean}

\medskip\noindent
{\bf Conference information.}

\href{%
https://sites.google.com/view/%
60thanniversary-dmat-puc-rio/Home
}{International Conference Celebrating
the 60th Anniversary of the Department
of Mathematics at PUC-Rio}

August 12--14, 2026,
RDC Auditorium, PUC-Rio,
Rio de Janeiro.

\medskip\noindent
{\bf Speaker.}

Paulo Orenstein, IMPA.

\medskip\noindent
{\bf Date and time.}

August 12, 2026, 11:15--12:15.

\medskip\noindent
{\bf Title.}

Robust and Efficient Mean Estimation
via Soft Shrinkage.

\medskip\noindent
{\bf Abstract.}

How should we weight observations when
estimating a mean? The sample mean is
efficient and minimax but fragile under
heavy tails; robust estimators resist
heavy tails but typically sacrifice
efficiency. We show a soft-shrinkage
calibration improves on both fronts:
observations are smoothly downweighted
by their distance from a pilot estimate,
with the shrinkage scale calibrated so
that the total rejected weight equals a
prescribed budget. Any reasonable pilot
is upgraded to sub-Gaussian concentration
with the minimax constant $\sqrt{2}$ under
finite variance alone, and the budget
doubles as an exact breakdown point. The
mechanism trades variance for bias, and
the bias can be quantified exactly. This
yields a second-order theory establishing
an MSE expansion that identifies when
shrinkage beats the sample mean,
Berry--Esseen and moderate deviation
approximations, and asymptotically exact
confidence intervals. Beyond the theory,
the results give concrete guidance for
designing mean estimators and bridge
modern sub-Gaussian estimation with
classical robust statistics. This is
joint work with Ant\^onio Cat\~ao.

\medskip\noindent
{\bf Notes.}
The talk starts introducing the mean
$\mu$ of a sample $X_i$, which is 
the sum of all the variables divided
by the size of the sample. This is
an efficient estimate, but not robust.
Robustness can be explained with atypical
data: suppose that you have a very atypical
datum in the sample. This will move the
mean $\mu$ quite a lot, but this isn't very
good since it was just one extraordinary
datum. Then we introduce the trimmed mean,
which uses a new variable to trim, or
remove, the tails of the data, i.e.
the most extraordinary data points
to avoid the case we just explained.
The trimmed mean is robust, but looses
efficiency since we need to chop off a 
piece of our data, which isn't ideal.

The talk aims at introducing a new 
estimate which lies in between the 
usual mean and the trimmed mean.
The idea is to introduce weights to every
data using a variable  $\alpha$ (which
depends on the whole data).
Such a new estimate is both robust
and efficient.

Developing these kinds of statistical
gadgets has applications in mathematics,
science, computer science, artificial
intelligence, etc. In particular, this can
be thought of as a more sophisticated method
than simple gradient descent used in machine
learning methods, the one used in the AI
course at EGAXII.


\bibliography{my}
\bibliographystyle{amsalpha}

\end{document}
